\documentclass[11pt,onecolumn]{article}

\newenvironment{cwpabstract}
{%
  \begin{list}{}{%
    \setlength{\leftmargin}{0.32\textwidth} % spostamento laterale
    \setlength{\rightmargin}{0pt}
    \setlength{\listparindent}{0pt}
    \setlength{\itemindent}{0pt}
    \setlength{\parsep}{0pt}
  }
  \item[]
  \noindent\textbf{ABSTRACT}\par
  \vspace{0.5em}
  \normalsize
}
{%
  \end{list}
}

\newenvironment{cwpkeywords}
{%
  \begin{list}{}{%
    \setlength{\leftmargin}{0.32\textwidth}
    \setlength{\rightmargin}{0pt}
  }
  \item[]\textbf{Key words: }
}
{%
  \end{list}
}

\usepackage{times}
\usepackage{xcolor}


\usepackage{geometry}
\geometry{
  paperwidth=8.5in,
  paperheight=11in,
  top=0.75in,
  bottom=1.0in,
  left=1.0in,
  right=1.0in,
  textwidth=6.5in,
  textheight=8.75in
}

\usepackage{setspace}
\setstretch{1.0}


\usepackage{titlesec}

\titleformat{\section}
  {\bfseries\normalsize}
  {\thesection}
  {1em}
  {\MakeUppercase}

\titleformat{\subsection}
  {\bfseries\normalsize}
  {\thesubsection}
  {1em}
  {}

\titleformat{\subsubsection}
  {\itshape\normalsize}
  {\thesubsubsection}
  {1em}
  {}

\usepackage{graphicx}
\usepackage{subcaption}
\usepackage{appendix}

% Packages used by default CWP reports built by M8R
\usepackage{natbib,amsmath,graphicx,color,amssymb,amsbsy,lineno,setspace,algorithm2e,lipsum,hyperref}



\title{Accoppiamento FEM/BEM: implementazione delle condizioni di dominio aperto nel codice FEM 2D
elettrostatico \thanks{Progetto realizzato nell'ambito del corso \textit{Computational Electromagnetics}}} 

\author{Matteo Morana \\
Master Student in Mathematical Engeneering at Politecnico di Milano \\ email: matteo.morana@mail.polimi.it}

\date{April 13, 2026}

\begin{document}


\maketitle


\begin{cwpabstract}
\lipsum[2]
\end{cwpabstract}


\begin{cwpkeywords}
Keyword1, Keyword2, Keyword3 
\end{cwpkeywords}


\section{Introduction}
A classic example of FEM/BEM coupling in an electrostatic problem is the computation of the electric field generated by a parallel-plate capacitor. In a simulation of this type, it is often useful to visualize not only the behavior of the field inside the capacitor, but also outside it. In Physics II courses, the external field is usually assumed to be zero; in reality, due to edge effects, it is not negligible, and in some applications it can play a significant role. This is the main motivation for analyzing in greater detail the behavior of the electric field in such a problem.

In this project, we consider the relatively simple case of a parallel-plate capacitor. Besides the charges present on the plates, no other sources are considered in the surrounding space. The differential equation to be solved is Maxwell's first law in the case with no density of charge inside the capacitor. ~\eqref{eq:prima_legge_maxwell_semplificata}, in the formulation of this problem we express the law by treating the electric field as the gradient of the potential ~\eqref{eq:prima_legge_maxwell_semplificata_potenziale}.

\begin{equation}
	\nabla \cdot E = 0
\label{eq:prima_legge_maxwell_semplificata}
\end{equation}

\begin{equation}
	\Delta V = 0
\label{eq:prima_legge_maxwell_semplificata_potenziale}
\end{equation}

A well-known first approach consists in using the Finite Element Method (FEM). One defines a domain large enough to contain the capacitor and imposes a zero potential on its boundary, assuming that at a sufficiently large distance the potential effectively tends to zero. However, this method has some drawbacks: it requires very large domains and therefore large meshes, resulting in increased computational time. To mitigate this issue, one may employ non-uniform meshes, refining the discretization only in regions where higher accuracy is required while keeping a coarser mesh elsewhere. This approach can significantly reduce the computational cost. However, the implementation of non-uniform meshes is not straightforward, as special care must be taken in handling the transitions between regions with different mesh resolutions. Moreover, it is necessary to identify the areas where refinement is most beneficial; in this problem, the most critical points are typically located near the edges and corners of the capacitor. Computing the external field with FEM therefore leads to significant computational waste, and in any case the potential cannot be represented up to infinity, introducing an unavoidable approximation.

The Boundary Element Method (BEM), on the other hand, offers a more efficient solution to the problem. By reducing the dimension of the domain by one order, it allows substantial computational advantages. By its nature, BEM enables the accurate computation of the unknown variable at any internal point of a domain, solving the problem only on its boundary. BEM is also applicable to very large or unbounded domains, provided that the asymptotic behavior of the potential is known a priori and that it converges as $1/r^2$. In such cases, BEM can correctly represent the behavior at infinity without introducing artificial approximations.

The goal of this work is therefore to combine FEM and BEM to obtain more accurate results at lower computational cost. The idea is to use FEM to determine the solution accurately inside the capacitor and employ BEM to extend it outward to infinity in an exact manner. The coupling mechanism between the two methods will be described in detail in the dedicated section. ~\cite{fem_bem_text}

\section{Coupling Strategy}
\subsection{FEM formulation}
In this section we derive the FEM formulation for our problem, we do not perform the FEM as usual, since we need to achieve a specific formulation in order to couple FEM and BEM.


The problem to solve is 

\begin{equation}
    \Delta u = 0 \quad \text{in } \Omega.
\label{eq:problema_laplaciano}
\end{equation}

Let's rewrite the problem in the weak formulation 

\begin{equation}
\begin{aligned}
&\text{Find } u \in H^1(\Omega) \quad \text{s.t}
&\int_{\Omega} \nabla u \cdot \nabla v \, dx  \, -\int_{\partial\Omega} \frac{\partial u}{dn}v\,ds = 0
\qquad \forall v \in H_{g}^1(\Omega) := \{v \in H^1  \text{ s.t. } v ... \}
\end{aligned}
\label{eq:weak_formulation}
\end{equation}

We did not write boundary conditions for the problem. this is because we have neither Dirichlet nor Neumann boundary conditions; we have a relation between $u$ and $\frac{\partial u}{dn}$ thanks to the BEM method. So, for now, let's write the problem in this way without imposing any bcs.

Now we consider that to have equation ~\eqref{eq:weak_formulation} verified $\forall v$ it's sufficient to verify for each element of the base of $H^1$. We also discretize the problem and so do not solve the problem for $v \in H^1$ but for $v_h \in V_h \subset H^1$. These consideration leads to $N_h$ equations, where $N_h$ is the dimension of the space $V_h$. Now we observe that $u \in V_h$ therefore we can write $u$ as $\sum_{j=1}^{N_h} u_jv_j$ where $v_j$ are the basis of $V_h$. We can now write the problem as in ~\eqref{eq:discrete_formulation} where $v_i$ forms the base of the $V_h$ space. We can see ~\eqref{eq:discrete_formulation} as a system of $N_h$ equations in $N_h$ unknowns which are $u_j$. If we achieve $u_h$ we can perform $\sum_{j=1}^{N_h} u_jv_j$ and obtain $u \in V_h$ which is the solution of our problem.

\begin{equation}
\begin{aligned}
&\text{Find } u_j \,,\, j = 1,...,N_h\quad \text{s.t}
&\int_{\Omega} \nabla (\sum_{j=1}^{N_h} u_jv_j) \cdot \nabla v_i \, dx  \, -\int_{\partial\Omega} \frac{\partial (\sum_{j=1}^{N_h} u_jv_j)}{dn}v_i\,ds = 0
\qquad \forall v_i \,,\, i = 1,...,N_h
\end{aligned}
\label{eq:discrete_formulation}
\end{equation}

With this notation we lose information about the normal derivative since we can rearrange the term and take out $u_j$ that are constant from the derivative. So to mantain a $\frac{\partial u}{dn}$ we consider that $\frac{\partial u}{dn}$ can be written as $\sum_{j=1}^{N_h} x_jv_j$ where $v_j$ are the basis of $V_h$ and $x_j$ are coefficient. In fact since we can select the space $V_h$, it's possibile to choose a space who garantee that both $u$ and $\frac{\partial u}{dn}$ are element of this space an therefore written as a proper linear combination of the basis of the space. We rewrite ~\eqref{eq:discrete_formulation} as ~\eqref{eq:discrete_formulation_2}, in the new formulation $x_i$ represents the coefficient to find to achieve the normal derivative. To highlight the meaning of $x_i$ we report also ~\eqref{eq:discrete_formulation_3} in which $x_i$ are called $\frac{\partial u_i}{dn}$, is important to pay attention to this notation because we are not appling the normal derivative operator to $u_i$ which is a coefficient, and the operator would not be meaningful since the operation would result in 0. We are just denoting the coefficient with the symbol $\frac{\partial u_i}{dn}$

\begin{equation}
\begin{aligned}
&\text{Find } u_j ,x_j \,,\, j = 1,...,N_h\quad \text{s.t}
&\int_{\Omega} \nabla (\sum_{j=1}^{N_h} u_jv_j) \cdot \nabla v_i \, dx  \, -\int_{\partial\Omega} (\sum_{j=1}^{N_h} {\color{red}x_j}v_j)v_i\,ds = 0
\qquad \forall v_i \,,\, i = 1,...,N_h
\end{aligned}
\label{eq:discrete_formulation_2}
\end{equation}


\begin{equation}
\begin{aligned}
&\text{Find } u_j, \frac{\partial u_j}{dn}  \,,\, j = 1,...,N_h\quad \text{s.t}
&\int_{\Omega} \nabla (\sum_{j=1}^{N_h} u_jv_j) \cdot \nabla v_i \, dx  \, -\int_{\partial\Omega} (\sum_{j=1}^{N_h} {\color{red}\frac{\partial u_j}{dn} }v_j)v_i\,ds = 0
\qquad \forall v_i \,,\, i = 1,...,N_h
\end{aligned}
\label{eq:discrete_formulation_3}
\end{equation}


\begin{equation}
\begin{aligned}
&\text{Find } u_j,\frac{\partial u_j}{dn}  \,,\, j = 1,...,N_h\quad \text{s.t}
&\int_{\Omega}\sum_{j=1}^{N_h} ({\color{blue} u_j}  {\color{green} \nabla v_j \cdot \nabla v_i})  \, dx  \, -\int_{\partial\Omega} (\sum_{j=1}^{N_h} {\color{red}\frac{\partial u_j}{dn} }{\color{green}v_jv_i})\,ds = 0
\qquad \forall v_i \,,\, i = 1,...,N_h
\end{aligned}
\end{equation}


\begin{equation}
\begin{aligned}
&\text{Find } u_j,\frac{\partial u_j}{dn}  \,,\, j = 1,...,N_h\quad \text{s.t}
&\sum_{j=1}^{N_h} ({\color{blue} u_j}  {\color{green}\int_{\Omega} \nabla v_j \cdot \nabla v_i\, dx})    \, - (\sum_{j=1}^{N_h} {\color{red}\frac{\partial u_j}{dn} }{\color{green}\int_{\partial\Omega}v_jv_i\,ds}) = 0
\qquad \forall v_i \,,\, i = 1,...,N_h
\end{aligned}
\end{equation}


Doing some transformation we can rewrite the problem as
\begin{equation}
\begin{aligned}
\text{Find } u,\frac{\partial u}{\partial n} \text{ s.t. } \quad
&{\color{green} S_i}\, {\color{blue} u} 
- 
{\color{green} T_i}\,
{\color{red} \frac{\partial u}{\partial n}} 
= 0 \qquad \forall v_i, \quad i = 1, \dots, N_h \quad \text{where:}
\end{aligned}
\label{eq:discrete_formulation_4}
\end{equation} 

\begin{equation}
A_i 
= 
\begin{bmatrix}
\int_{\Omega}\nabla v_1 \cdot \nabla v_i &
\int_{\Omega}\nabla v_2 \cdot \nabla v_i &
\dots &
\int_{\Omega}\nabla v_{N_h} \cdot \nabla v_i
\end{bmatrix}
\end{equation}


\begin{equation}
\mathbf{u} = 
\begin{bmatrix}
u_1 & u_2 & \dots & u_{N_h}
\end{bmatrix}^T,
\end{equation}

\begin{equation}
T_i =
\begin{bmatrix}
\int_{\partial\Omega}v_1 v_i &
\int_{\partial\Omega}v_2 v_i &
\dots &
\int_{\partial\Omega}v_{N_h} v_i &
\end{bmatrix}
\end{equation}



\begin{equation}
\frac{\partial \mathbf{u}}{\partial n} =
\begin{bmatrix}
\frac{\partial u_1}{\partial n} &
\frac{\partial u_2}{\partial n} &
\dots &
\frac{\partial u_{N_h}}{\partial n}
\end{bmatrix}^T.
\end{equation}


Visualize the $N_h$ equation as a system we can write

\begin{equation}
\begin{aligned}
&\text{Find } u,\frac{\partial u}{\partial n} \quad \text{s.t}
&{\color{green}S}{\color{blue}u}  \, - {\color{green}T}{\color{red}\frac{\partial u}{dn}} = 0
\end{aligned}
\label{eq:discrete_formulation_5}
\end{equation}
%% ------------------------------------------------------------

\subsection{BEM formulation}
In this section we derive the BEM formulation for our problem. The key idea behind BEM is to reformulate the problem defined on $\Omega$ as an equivalent problem defined only on the boundary $\partial\Omega$, reducing the dimensionality of the problem.

The starting point is the fundamental solution of the Laplace operator. We define $G(x,y)$ as the function satisfying

\begin{equation}
    \Delta_x G(x,y) = \delta(x-y) \quad \text{in } \mathbb{R}^d,
\label{eq:fundamental_solution}
\end{equation}

where $\delta$ is the Dirac delta distribution and $d$ is the space dimension. The explicit expression of $G$ is:

\begin{equation}
G(x,y) = 
\begin{cases}
\dfrac{1}{2\pi} \log |x - y| & \text{if } d = 2, \\[6pt]
\dfrac{1}{4\pi |x - y|} & \text{if } d = 3.
\end{cases}
\label{eq:green_function}
\end{equation}

\subsubsection{Green's third identity}
Let $u \in H^1(\Omega)$ be a solution of~\eqref{eq:problema_laplaciano}. Applying Green's second identity to $u$ and $G(\cdot, y)$ for a {\color{purple}fixed $y \in \Omega$ } yields the \textit{representation formula}:

\begin{equation}
u(y) = \int_{\partial\Omega} G(x,y)\, \frac{\partial u}{\partial n}(x)\, ds_x 
\,-\, \int_{\partial\Omega} \frac{\partial G}{\partial n_x}(x,y)\, u(x)\, ds_x, 
\qquad \forall {\color{purple} y \, \text{(fixed)}} \in \Omega.
\label{eq:representation_formula}
\end{equation}

This formula expresses $u$ at any interior point $y$ solely in terms of the boundary values of $u$ and $\frac{\partial u}{\partial n}$. We emphasize that, since $(y)$ is fixed when evaluating the representation formula, the integration is performed with respect to the boundary variable $(x)$; hence the boundary element is written ($ds_{x}$), indicating that the integration takes place over the boundary while $(y)$ is kept fixed. We define the two integral operators appearing in~\eqref{eq:representation_formula}:

\begin{equation}
(\mathcal{V}\phi)(y) := \int_{\partial\Omega} G(x,y)\, \phi(x)\, ds_x,
\label{eq:single_layer}
\end{equation}

\begin{equation}
(\mathcal{K}\phi)(y) := \int_{\partial\Omega} \frac{\partial G}{\partial n_x}(x,y)\, \phi(x)\, ds_x,
\label{eq:double_layer}
\end{equation}

where $\mathcal{V}$ is called the \textit{single layer operator} and $\mathcal{K}$ is called the \textit{double layer operator}. With this notation,~\eqref{eq:representation_formula} becomes:

\begin{equation}
u(y) = \mathcal{V}\left(\frac{\partial u}{\partial n}\right)(y) - \mathcal{K}(u)(y), \qquad \forall y \in \Omega.
\end{equation}

\subsubsection{Boundary Integral Equation}
To obtain a problem posed exclusively on $\partial\Omega$, we take the trace of the representation formula, i.e., we let $y \to \partial\Omega$. Due to the singularity of $G$, this limit must be handled carefully and leads to the \textit{boundary integral equation} (BIE):

\begin{equation}
\frac{1}{2}\,u(y) = \mathcal{V}\left({\color{red}\frac{\partial u}{\partial n}}\right)(y) - \mathcal{K}({\color{blue}u})(y), 
\qquad{\color{purple} \forall y \in \partial\Omega}.
\label{eq:bie}
\end{equation}

The factor $\frac{1}{2}$ arises from the jump relations of the double layer potential across a \textbf{smooth} boundary. Equation~\eqref{eq:bie} is a relation between $u|_{\partial\Omega}$ and $\frac{\partial u}{\partial n}|_{\partial\Omega}$: given one of the two, one can in principle recover the other. \textit{Remember that the factor $\frac{1}{2}$ is in general a value $\alpha$, depends on $\alpha$ the equation is valid for inside the domain, on a smooth border, on a border with corners or outside of the domain. In general $\alpha = \frac{\theta}{2\pi}$}. \textbf{Since the boundary $\partial\Omega$ is a rectangle, the numerical implementation assigns $\alpha = \frac{1}{4}$ to the four corner nodes and $\alpha = \frac{1}{2}$ to the remaining smooth boundary nodes. For simplicity of notation, we write $\frac{1}{2}$ throughout the paper.}

\subsubsection{Discretization}
We now discretize~\eqref{eq:bie}. Let $V_h \subset H^{1/2}(\partial\Omega)$ be a finite-dimensional space of dimension $N_h^{\partial}$, with basis $\{\psi_j\}_{j=1}^{N_h^{\partial}}$. We approximate the boundary unknowns as:

\begin{equation}
{\color{blue}u_h = \sum_{j=1}^{N_h^{\partial}} u_j\, \psi_j}, 
\qquad 
{\color{red}\frac{\partial u_h}{\partial n} = \sum_{j=1}^{N_h^{\partial}} \frac{\partial u_j}{\partial n}\, \psi_j}.
\label{eq:bem_approximation}
\end{equation}

Following the same notation introduced for FEM, \textbf{we remark that $\frac{\partial u_j}{\partial n}$ denotes a coefficient}, not the application of the normal derivative operator to the scalar $u_j$. Substituting~\eqref{eq:bem_approximation} into~\eqref{eq:bie} and testing against $\psi_i$. Below the steps to obtain the discrete system for the BEM formulation ~\eqref{eq:bem_discrete}. 

\begin{equation}
\frac{1}{2}\,({\color{blue}\sum_{j=1}^{N_h^{\partial}} u_j\, \psi_j})(y) = \mathcal{V}\left({\color{red}\sum_{j=1}^{N_h^{\partial}} \frac{\partial u_j}{\partial n}\, \psi_j}\right)(y) - \mathcal{K}({\color{blue}\sum_{j=1}^{N_h^{\partial}} u_j\, \psi_j})(y), 
\qquad{\color{purple} \forall y \in \partial\Omega}.
\label{eq:bem_discrete_step1}
\end{equation}

\begin{equation}
\int_{\partial\Omega}\frac{1}{2}\,({\color{blue}\sum_{j=1}^{N_h^{\partial}} u_j\, \psi_j})\psi_i ds= \int_{\partial\Omega}\mathcal{V}\left({\color{red}\sum_{j=1}^{N_h^{\partial}} \frac{\partial u_j}{\partial n}\, \psi_j}\right)(y)\psi_i \, ds_y - \int_{\partial\Omega}\mathcal{K}({\color{blue}\sum_{j=1}^{N_h^{\partial}} u_j\, \psi_j})(y)\psi_i \, ds_y, 
\qquad{\color{purple} \forall y \in \partial\Omega} \quad \forall i \in N_h^{\partial}.
\label{eq:bem_discrete_step2}
\end{equation}

\begin{equation}
\frac{1}{2}\sum_{j=1}^{N_h^{\partial}} {\color{blue}u_j} \int_{\partial\Omega} \psi_j \psi_i\, ds
=
\sum_{j=1}^{N_h^{\partial}} {\color{red}\frac{\partial u_j}{\partial n} }\int_{\partial\Omega} \mathcal{V}(\psi_j)\, \psi_i\, ds_y
\,-\,
\sum_{j=1}^{N_h^{\partial}} {\color{blue}u_j} \int_{\partial\Omega} \mathcal{K}(\psi_j)\, \psi_i\, ds_y,
\qquad \forall i = 1,\dots,N_h^{\partial}.
\label{eq:bem_discrete}
\end{equation}

We define the matrices:

\begin{equation}
M_{ij} = \int_{\partial\Omega} \psi_j\, \psi_i\, ds,
\end{equation}

\begin{equation}
G_{ij} = \int_{\partial\Omega} \mathcal{V}(\psi_j)\, \psi_i\, ds,
\end{equation}

\begin{equation}
K_{ij} = \int_{\partial\Omega} \mathcal{K}(\psi_j)\, \psi_i\, ds.
\end{equation}

With this notation,~\eqref{eq:bem_discrete} can be written compactly as:

\begin{equation}
\begin{aligned}
&\text{Find } u,\frac{\partial u}{\partial n} \quad \text{s.t.}
&\left(\frac{1}{2}M + K\right) {\color{blue} u} = G\, {\color{red}\frac{\partial u}{\partial n}}.
\end{aligned}
\label{eq:bem_system_general}
\end{equation}

The problem ~\eqref{eq:bem_system_general} can be written as:

\begin{equation}
\begin{aligned}
&\text{Find } u,\frac{\partial u}{\partial n} \quad \text{s.t.}
&H{\color{blue} u} = G\, {\color{red}\frac{\partial u}{\partial n}}.
\end{aligned}
\label{eq:bem_system}
\end{equation}

\textbf{We remark that problem ~\eqref{eq:bem_system} is valid to find $u,\frac{\partial u}{\partial n} $ on the border of the domain $\Omega$}

\subsubsection{G is symmetric}
\begin{equation}
G_{ij} = \int_{\partial\Omega} \left[ \int_{\partial\Omega} G(x,y)\, 
\psi_j\, ds_x \right] \psi_ids_y
\label{eq:bem_G_sym_step_1}
\end{equation}

\begin{equation}
G_{ij} = \int_{\partial\Omega} \int_{\partial\Omega} G(x,y)\, 
\psi_j\, \psi_ids
\label{eq:bem_G_sym_step_2}
\end{equation}

therefore $G_{ij} = G_{ji}$

\subsection{FEM-BEM Coupling}

Before combining the two formulations, it is useful to clarify the structure of the discrete FEM system and the number of unknowns it involves.  
The FEM system~\eqref{eq:discrete_formulation_5} provides $N_h$ equations. However, in the above formulation both the nodal values of $u$ and the normal derivative $\frac{\partial u}{\partial n}$ appear as unknowns. If no further structure is exploited, the system would therefore involve $2N_h$ unknowns.

To reduce the number of unknowns and make the coupling feasible, additional structure can be imposed through the choice of the basis functions $v_i$.

If Lagrange nodal basis functions are used, they satisfy $v_i(x_j)=\delta_{ij}$. As a consequence, for any interior node $i$ the corresponding basis function vanishes at all boundary nodes, and therefore $v_i|_{\partial\Omega}=0$. In this case the boundary integral vanishes whenever at least one of the basis function is associated with interior nodes. In particular,
\[
T_{ij} = \int_{\partial\Omega} v_i \, v_j = 0 
\qquad \text{if } i \text{ or } j \text{ corresponding to an interior nodes}.
\]

\textbf{We remark that this is \textit{not} a general property of $H^1$ bases, but rather a consequence of a specific choice of basis.} 

This property implies that the coefficients of $\frac{\partial u}{\partial n}$ associated with interior nodes do not appear in the system. Only the coefficients corresponding to nodes located on the boundary remain. If we denote by $N_h^{\partial}$ the number of boundary nodes, the set of unknowns reduces to the $N_h$ nodal values of $u$ together with the $N_h^{\partial}$ boundary coefficients of $\frac{\partial u}{\partial n}$. The FEM formulation therefore provides $N_h$ equations in $N_h + N_h^{\partial}$ unknowns.

For an \textit{interior} node $i$, since $v_i|_{\partial\Omega}=0$, the boundary contribution vanishes and the corresponding equation reduces to
\begin{equation}
S_i\, u = 0,
\label{eq:interior_row}
\end{equation}
which is already closed in the unknowns $u$.

For a \textit{boundary} node $i$, instead, the boundary term is generally nonzero and the equation reads
\begin{equation}
S_i\, u - T_i\, \frac{\partial u}{\partial n} = 0,
\label{eq:boundary_row}
\end{equation}
so that the additional unknown $\frac{\partial u}{\partial n}$ appears explicitly. These $N_h^{\partial}$ equations remain open.

\textbf{We remark that $\frac{\partial u}{\partial n}$ is not zero in the interior of the domain. However, within the present formulation we do not have access to its values at interior nodes. These quantities do not enter the discrete system and therefore cannot be determined at this stage. They can only be reconstructed \textit{a posteriori}, once the system has been solved.}
This is precisely where the BEM formulation will provide the missing relations.
The BEM system~\eqref{eq:bem_system} provides exactly $N_h^\partial$ equations, but crucially only involves quantities defined on $\partial\Omega$: it is a relation between $u|_{\partial\Omega}$ and $\frac{\partial u}{\partial n}|_{\partial\Omega}$. The coupling strategy is therefore the following: \textit{for boundary nodes, replace the open FEM equations~\eqref{eq:boundary_row} with the BEM equations~\eqref{eq:bem_system}}. This is conceptually analogous to imposing Dirichlet conditions in a standard FEM problem: in that case some rows of the system matrix are replaced by $u_j = g_j$; here, the BEM equation plays the same role for the boundary rows, providing the information that the FEM formulation alone cannot supply.

After the substitution, the total system has $N_h + N_h^\partial$ equations in $N_h + N_h^\partial$ unknowns and takes the form~\eqref{eq:coupled_system}:

\begin{equation}
\begin{bmatrix}
S_{II} & S_{IB} & 0 \\
S_{BI} & S_{BB} & -T_{BB} \\
0 & H & -G
\end{bmatrix}
\begin{bmatrix}
u_I \\ u_B \\ \dfrac{\partial u_B}{\partial n}
\end{bmatrix}
=
\begin{bmatrix}
0 \\ 0 \\ 0
\end{bmatrix},
\label{eq:coupled_system}
\end{equation}

where the subscripts $I$ and $B$ denote interior and boundary degrees of freedom respectively, $S_{II}$ and $S_{IB}$, $S_{BI}$, $S_{BB}$ are the blocks of the FEM stiffness matrix $S$ partitioned accordingly, $T_{BB}$ is the boundary mass block from the FEM formulation, and $M$, $K$, $V$ are the BEM matrices defined in~\eqref{eq:bem_system}. 

The block structure of~\eqref{eq:coupled_system} makes the role of each component transparent: the first block of rows, coming entirely from FEM, governs the interior of the domain and does not involve $\frac{\partial u}{\partial n}$; the second block of rows, also from FEM, couples interior and boundary values of $u$ with the normal derivative; the third block of rows, from BEM, closes the system by providing the missing relation between $u_B$ and $\frac{\partial u_B}{\partial n}$ through the integral operators on $\partial\Omega$.


\subsubsection{Remarks on the block structure}
The blocks $S_{IB}$ and $S_{BI}$, are subblocks of the FEM stiffness matrix, whose entries are $S_{ij} = \int_{\Omega} \nabla v_i \cdot \nabla v_j \, dx$. Since this integral is over the \textit{whole domain} $\Omega$ and not just the boundary, it is in general nonzero even when one index is interior and the other is on the boundary: if the supports of $v_i$ and $v_j$ overlap near $\partial\Omega$, the integral is nonzero. Therefore $S_{IB} \neq 0$ and $S_{BI} \neq 0$ in general.

\subsubsection{Remarks on the BEM normal direction}

Since we aim to use BEM to recover information outside the domain, the outward-pointing normal with respect to the exterior region must be considered. In the classical formulation, this corresponds to using $-n$ in place of the standard outward normal. By factoring out the negative sign, the BEM formulation and its numerical implementation remain unchanged, yielding

\begin{equation}
    \frac{\partial u_B}{\partial{(-n)}} = G^{-1}H u_B.
    \label{eq:elimination}
\end{equation}

Pulling the sign outside the normal derivative introduces a sign change in the linear system.

The BEM can be applied to unbounded exterior domains provided that the quantity of interest — in this case the potential $u$ — decays at infinity at least as $1/r^2$. This decay condition is essential to guarantee the integrability of the single layer operator ~\eqref{eq:single_layer} and the well-posedness of the boundary integral formulation. In our setting, the capacitor can be regarded as a dipole when viewed from sufficiently large distances. The potential of a dipole decays as $1/r^2$ \cite{potential_dipole}, which ensures that the decay condition is satisfied and the BEM formulation on the exterior domain is well-posed.


\subsubsection{Reduced system after elimination}

We can proceed by analogy of imposing dirichlet condition and go for elimination, \textbf{we assume $G$ is invertible}:

\begin{equation}
\frac{\partial u_B}{\partial n} = -G^{-1}H u_B.
\label{eq:elimination}
\end{equation}

Substituting~\eqref{eq:elimination} into the second block row of~\eqref{eq:coupled_system}:

\begin{equation}
S_{BI}\, u_I + S_{BB}\, u_B + T_{BB}\, G^{-1}H u_B = 0,
\end{equation}

In this way we eliminate the third block row and the reduced system have only the unknowns $u_I$ and $u_B$:

\begin{equation}
\begin{bmatrix}
S_{II} & S_{IB} \\
S_{BI} & S_{BB} + T_{BB}\, G^{-1}H
\end{bmatrix}
\begin{bmatrix}
u_I \\ u_B
\end{bmatrix}
=
\begin{bmatrix}
0 \\ 0
\end{bmatrix}.
\label{eq:reduced_system}
\end{equation}

The matrix in~\eqref{eq:reduced_system} has size $N_h \times N_h$ and is the sum of the standard FEM stiffness contribution and a correction term $T_{BB}\, G^{-1}H$ that is dense and restricted to boundary degrees of freedom, reflecting the nonlocal nature of the BEM operator. Once $u_I$ and $u_B$ are obtained by solving~\eqref{eq:reduced_system}, the normal derivative can be recovered a posteriori via~\eqref{eq:elimination}.
 In the MATLAB implementation of this method, the product $G^{-1}H$ is computed using the backslash operator \texttt{\textbackslash}, i.e. by solving the linear system $GX = H$.


\subsubsection{Matrix entries}

For convenience of implementation, we collect here the explicit expression of the generic $(i,j)$ entry of each matrix appearing in~\eqref{eq:reduced_system}.

The FEM stiffness matrix $S$ has entries:
\begin{equation}
S_{ij} = \int_{\Omega} \nabla v_i \cdot \nabla v_j \, dx,
\end{equation}
where $v_i, v_j \in V_h$ range over all $N_h$ basis functions of the domain discretization.

The FEM boundary mass matrix $T$ has entries:
\begin{equation}
T_{ij} = \int_{\partial\Omega} v_i\, v_j \, ds,
\end{equation}
where in the block notation of~\eqref{eq:coupled_system} only the subblock $T_{BB}$ is relevant, with $i,j$ ranging over boundary nodes only.

The BEM single layer matrix $G$ has entries:
\begin{equation}
G_{ij} = \int_{\partial\Omega} \mathcal{V}(\psi_j)(y)\, \psi_i(y) \, ds_y = \int_{\partial\Omega}\int_{\partial\Omega} G(x,y)\, \psi_j(x)\, \psi_i(y) \, ds_x \, ds_y,
\end{equation}
where $\psi_i, \psi_j \in V_h^\partial$ are the basis functions of the boundary discretization and $G(x,y)$ is the fundamental solution~\eqref{eq:green_function}. Note that $G_{ij}$ involves a double boundary integral and is symmetric.

The BEM double layer matrix $H$ has entries:
\begin{equation}
H_{ij} = \frac{1}{2}\int_{\partial\Omega} \psi_i\, \psi_j \, ds + \int_{\partial\Omega} \mathcal{K}(\psi_j)(y)\, \psi_i(y) \, ds_y = \frac{1}{2}M_{ij} + \int_{\partial\Omega}\int_{\partial\Omega} \frac{\partial G}{\partial n_x}(x,y)\, \psi_j(x)\, \psi_i(y) \, ds_x \, ds_y,
\end{equation}
where the first term is the scaled boundary mass contribution and the second involves the normal derivative of the fundamental solution with respect to the source point $x$. Unlike $G$, the matrix $H$ is in general non-symmetric due to the directional nature of the double layer kernel.

\subsubsection{About the Code Implementation}

The implementation is based on two pre-existing solvers: a FEM code \cite{fem_code} and a BEM code \cite{bem_code}. Since the two codes are independent, the node ordering conventions they adopt differ: the BEM matrices $\mathbf{G}$ and $\mathbf{H}$ are assembled with boundary nodes ordered counterclockwise (CCW), while the FEM system uses an arbitrary node ordering. To interface the two solvers with minimal modifications to their original implementations, a set of auxiliary functions was developed.

\medskip\noindent
In particular, to align the two systems, a permutation vector \texttt{perm} is computed via the MATLAB command \texttt{[\textasciitilde, perm] = ismember(B\_nodes, B\_ordered)}, where \texttt{B\_nodes} contains the FEM node indices and \texttt{B\_ordered} contains the same indices sorted in CCW geometric order. The entry \texttt{perm(i)} returns the position of \texttt{B\_nodes(i)} within \texttt{B\_ordered}, so that the reordered matrices $\mathbf{G}(\texttt{perm},\texttt{perm})$ and $\mathbf{H}(\texttt{perm},\texttt{perm})$ are consistent with the FEM degree-of-freedom ordering.

\medskip\noindent
Throughout this paper, the permittivity constants $\varepsilon_0$ and $\varepsilon_r$ have not been explicitly discussed. This is because the relative permittivity is assumed to be $\varepsilon_r = 1$, and the vacuum permittivity $\varepsilon_0$ does not affect the solution as long as no source terms are present, i.e.\ as long as the problem remains a Laplace equation. The code, however, implements both $\varepsilon_0$ and $\varepsilon_r$ in full generality, so that it can be readily extended to more complex material configurations or Poisson-type problems.
%% ------------------------------------------------------------
\section{Using FEM-BEM coupled in the problem}
\subsection{Application to the parallel plate capacitor}

\subsubsection{Geometry and problem setup}

We consider two infinitely thin horizontal plates in 2D, placed at $y = +d/2$ and $y = -d/2$, 
extending from $x = -L/2$ to $x = +L/2$. The problem is $\Delta u = 0$ in $\mathbb{R}^2 
\setminus \{\text{plates}\}$ and Dirichlet conditions:
\begin{equation}
u = V_1 \quad \text{on the upper plate}, \qquad u = V_2 \quad \text{on the lower plate}.
\end{equation}

We introduce an artificial closed boundary $\Gamma$ (e.g. a rectangle containing both plates) 
that splits the domain into:
\begin{itemize}
    \item $\Omega_{\text{int}}$: the bounded interior region, treated with FEM,
    \item $\Omega_{\text{ext}} = \mathbb{R}^2 \setminus \overline{\Omega}_{\text{int}}$: 
    the unbounded exterior region, treated with BEM.
\end{itemize}

\subsubsection{Node partitioning}

We partition the degrees of freedom into three groups:
\begin{itemize}
    \item $u_D$: nodes on the plates, where the potential is \textbf{known}: $u_D = V_1$ or $u_D = V_2$,
    \item $u_F$: free interior nodes, where the potential is \textbf{unknown},
    \item $u_B$: nodes on $\Gamma$, where the potential is \textbf{unknown} but constrained by the BEM.
\end{itemize}

The union of nodes F and nodes D are the nodes I.

\subsubsection{Full block system}

The FEM formulation~\eqref{eq:discrete_formulation_5} written in block form over the three groups is:

\begin{equation}
\begin{bmatrix}
S_{FF} & S_{FD} & S_{FB} \\
S_{DF} & S_{DD} & S_{DB} \\
S_{BF} & S_{BD} & S_{BB}
\end{bmatrix}
\begin{bmatrix}
u_F \\ u_D \\ u_B
\end{bmatrix}
-
\begin{bmatrix}
0 \\ 0 \\ -T_{BB}\,G^{-1}H
\end{bmatrix}
= \begin{bmatrix}
0 \\ 0 \\ 0
\end{bmatrix}
\label{eq:full_block_system}
\end{equation}

where $T_{BB}$ acts only on $\Gamma$-nodes because interior basis functions vanish on $\Gamma$, 
and the plates lie inside $\Omega_{\text{int}}$ and therefore have no contribution to $T$. 

\subsubsection{Reduction of the system}

The row $D$ of~\eqref{eq:full_block_system} reads:
\begin{equation}
S_{DF}\, u_F + S_{DD}\, u_D + S_{DB}\, u_B = 0.
\label{eq:row_D}
\end{equation}
Since $u_D$ is already known by the Dirichlet condition, this equation is \textbf{not used to 
solve the system}. It will be recovered a posteriori to compute the surface charge density on 
the plates (see Section~\ref{sec:charge_density}).

We retain only rows $F$ and $B$. Since $u_D$ is known, we move it to the right-hand side. 
The reduced system is:

\begin{equation}
\begin{bmatrix}
S_{FF} & S_{FB} \\
S_{BF} & S_{BB} + T_{BB}\,G^{-1}H
\end{bmatrix}
\begin{bmatrix}
u_F \\ u_B
\end{bmatrix}
=
-\begin{bmatrix}
S_{FD} \\ S_{BD}
\end{bmatrix}
u_D,
\label{eq:reduced_system_capacitor}
\end{equation}

where $u_D$ is the vector of known Dirichlet values, with entries equal to $V_1$ or $V_2$ 
depending on which plate the node belongs to.

\subsubsection{Recovery of the full solution}\label{subsec:recovery_solution}

Once~\eqref{eq:reduced_system_capacitor} is solved, the complete solution is recovered in the 
following steps.

\textbf{Potential in $\Omega_{\text{int}}$.}
The FEM solution is fully determined: $u_F$ from the solved system, $u_D$ from the Dirichlet 
data, and $u_B$ from the solved system. The potential at any interior point is:
\begin{equation}
u_h(x) = \sum_j u_j\, v_j(x), \qquad x \in \Omega_{\text{int}}.
\end{equation}

\textbf{Normal derivative on $\Gamma$.}
The outward normal derivative on $\Gamma$ is recovered from the BEM 
relation:
\begin{equation}
\frac{\partial u_B}{\partial n} = -G^{-1} H\, u_B.
\end{equation}

\textbf{Potential in $\Omega_{\text{ext}}$.}
With $u_B$ and $\frac{\partial u_B}{\partial n}$ known on $\Gamma$, the potential at any 
exterior point $y \in \Omega_{\text{ext}}$ is recovered via the BEM representation 
formula~\eqref{eq:representation_formula}:
\begin{equation}
u(y) = \int_\Gamma G(x,y)\,\frac{\partial u_B}{\partial n}(x)\,ds_x 
\,-\, \int_\Gamma \frac{\partial G}{\partial n_x}(x,y)\,u_B(x)\,ds_x,
\qquad y \in \Omega_{\text{ext}},
\end{equation}
which in discrete form is a matrix-vector product using the same BEM matrices $G$ and $H$.

\subsubsection{Other phisical quantity}\label{subsec:compute_phisical_quantity}

\textbf{Surface charge density on the plates.}
\label{sec:charge_density}

The surface charge density on the capacitor plates is computed from the electric field distribution obtained from the solution. On a perfect conductor, the surface charge density is related to the normal component of the electric field by
\begin{equation}
\sigma = \varepsilon_0 E_n = -\varepsilon_0 \frac{\partial u}{\partial n}\bigg|_{\text{near plates}},
\end{equation}
where the normal derivative is evaluated in the dielectric region adjacent to the plate surface.

In the discrete FEM setting, we compute $\sigma$ by evaluating the gradient of the numerical solution in the elements adjacent to each plate. For each element $K$ sharing an edge (in 2D) with a plate node, we defined these elements as $K_{top}$ and $K_{bottom}$, we do the following (Case Top Plate):

\begin{enumerate}
    \item Compute the gradient $\nabla u|_{K_{top}}$ using the FEM shape functions and the nodal values $u_F$, $u_D$, and $u_B$.
    
    \item Evaluate the electric field:
    \[
    \mathbf{E}|_{K_{top}} = -\nabla u|_{K_{top}} .
    \]
    
    \item Project onto the outward normal $\mathbf{n}$ to the plate:
    \[
    E_n = \mathbf{E}|_{K_{top}} \cdot \mathbf{n}.
    \]
    
    \item Compute the local surface charge density:
    \[
    \sigma|_{K_{top}} = \varepsilon_0 E_n .
    \]
\end{enumerate}

\textbf{Capacitance}
\label{sec:capacitance}

The capacitance per unit length is computed using two independent methods, 
which serve as a mutual validation of the numerical solution.
The first method, known as the \textit{charge method}, computes the capacitance 
from the total charge stored on the plates as
\begin{equation}
    C = \frac{|Q|}{V_1 - V_2},
\end{equation}
where $Q$ denotes the total charge per unit length on one of the conductors, obtained 
by integrating the normal electric flux density over the corresponding plate. In the 
electrostatic formulation, this quantity is related to the normal derivative of the 
potential through Gauss's law
\begin{equation}
    Q = \varepsilon_0 \oint_{\Gamma_p} \frac{\partial u}{\partial n} \, d\Gamma,
\end{equation}
where $\Gamma_p$ denotes the plate surface.
The second method, known as the \textit{energy method}, estimates the capacitance 
from the electrostatic energy stored in the system. This approach is physically motivated 
by the fact that, for a real finite-length capacitor, the electric field is not confined 
to the gap between the plates: fringing fields extend beyond the edges and carry a 
non-negligible amount of energy. The electrostatic energy stored in the domain $\Omega$ 
is given by
\begin{equation}
    W = \frac{1}{2} \int_{\Omega} \varepsilon \, |\nabla u|^2 \, d\Omega,
\end{equation}
which, after discretization over the finite element mesh, becomes
\begin{equation}
    W \approx \frac{1}{2} \sum_{e} \varepsilon^{(e)} \left( E_x^{(e)2} + E_y^{(e)2} 
    \right) A^{(e)},
\end{equation}
where $\varepsilon^{(e)}$, $E_x^{(e)}$, $E_y^{(e)}$ and $A^{(e)}$ are respectively 
the permittivity, the two components of the electric field, and the area of the $e$-th 
element. The capacitance per unit length is then recovered as
\begin{equation}
    C = \frac{2W}{\Delta V^2},
\end{equation}
where $\Delta V = V_1 - V_2$ is the potential difference between the plates. Summing 
over the entire FEM domain ensures that the fringing field contribution is fully 
accounted for.
%% ------------------------------------------------------------
\section{Numerical Experiments}

To validate the FEM-BEM coupled model, we consider a two-dimensional parallel-plate capacitor 
as a benchmark problem. The geometry consists of two infinitely thin conducting plates of length 
$L_x$ separated by a gap $L_y$, embedded in free space ($\varepsilon_r = 1$). Dirichlet boundary 
conditions $u = V_1$ and $u = V_2$ are imposed on the upper and lower plates respectively, while the artificial coupling boundary $\Gamma$ is placed at a distance of approximately $10\,L_y$ 
from the plates on all sides. This choice is motivated by the spatial extent of the fringing 
fields: the field lines that curve around the plate edges decay on a length scale of order 
$L_y$ ~\cite{fringing_effect}, so a margin of $10\,L_y$ ensures that most of the fringing region is contained within 
$\Omega_{\text{int}}$ and adequately resolved by the FEM mesh. Moreover, since the bulk of the potential is concentrated in this region, applying FEM here guarantees higher accuracy where it is most needed, without any significant impact on the overall computational cost.
Placing $\Gamma$ closer to 
the plates would force the FEM to resolve steep field gradients on a coarse mesh, degrading 
accuracy; placing it farther would increase the size of $\Omega_{\text{int}}$ and the 
associated computational cost without a corresponding gain in accuracy, since the BEM handles 
the exterior domain exactly regardless of its extent.

A key requirement for a meaningful comparison with the analytical solution is that the aspect 
ratio $L_x/L_y$ is sufficiently large. For an ideal parallel-plate capacitor with infinite plates, 
the capacitance per unit length is given by the closed-form expression
\begin{equation}
    C_{\text{ideal}} = \varepsilon_0 \varepsilon_r \frac{L_x}{L_y},
    \label{eq:C_ideal}
\end{equation}
which neglects fringing fields at the plate edges. When $L_x/L_y \geq 10$, the fringing 
contribution is small (typically 5--10\%) and the numerical result can be meaningfully compared 
against~\eqref{eq:C_ideal}. For this reason we adopt $L_x = 0.10$ m, $L_y = 0.01$ m, giving 
$L_x/L_y = 10$ and $C_{\text{ideal}} = 88.5$ pF/m. The computational domain is set to 
$X_{mesh} = L_x + 2 \times 10\,L_y$ and $Y_{mesh} = L_y + 2 \times 10\,L_y$, and the mesh step is 
$\delta = 0.005$ m, yielding approximately 500 nodes and 90 nodes on $\Gamma$ or $\delta = 0.002$ m, yielding approximately 2900 nodes and 234 nodes on $\Gamma$. Simulations performed with $\delta = 0.001$ m did not yield any significant improvement over those obtained with $\delta = 0.002$ m; the additional computational cost is therefore not justified.

Several quantities are monitored to assess the correctness of the coupled solution as presented in section ~\ref{subsec:recovery_solution}.

\textbf{Electric potential $u$.} 
The potential must be monotone between the plates and lie strictly in the interval $[V_2, V_1]$ 
at every interior node. Outside the plates it should decay smoothly toward $\Gamma$, with no 
artificial oscillations near the coupling boundary. A contour plot of $u$ is used to verify 
these properties visually.

\textbf{Electric field $\mathbf{E} = -\nabla u$.}
Between the plates, and away from the edges, the field should be nearly uniform and directed 
perpendicularly to the plates, with magnitude close to $(V_1 - V_2)/L_y$. Near the plate edges 
a field concentration is physically expected — the field intensity diverges at the tips of 
infinitely thin conductors — so a small number of elements with anomalously large $|\mathbf{E}|$ 
is not a sign of error but an inherent feature of the geometry. (The electric field is computed using a function provided by the FEM code.)

\textbf{Normal derivative on $\Gamma$.}
The quantity $q = \frac{\partial u}{\partial n}$ on $\Gamma$, recovered as $q = G^{-1} H u_B$, represents the flux \textit{entering} $\Omega_{\mathrm{int}}$ through the artificial boundary. 
Since no sources are present in $\Omega_{\mathrm{int}}$, the total flux through 
$\Gamma$ must satisfy Gauss's law: the net integral of $q$ over $\Gamma$ should be zero 
(corresponding to equal and opposite charges on the two plates). 
A non-zero net flux would indicate an inconsistency in the coupling.

\textbf{Surface charge density on the plates.}
\label{sec:charge_density_results}

The procedure described in Section~\ref{sec:charge_density} is applied to 
each plate of the capacitor. A few practical remarks are worth noting.

For linear finite elements, the gradient $\nabla u|_K$ is constant within 
each element, so the charge density $\sigma_i$ is piecewise constant along 
the plate. As a consequence, the resolution of the charge distribution depends 
directly on the mesh refinement along the plates.

In the case of a parallel-plate capacitor with a uniform gap and a homogeneous 
dielectric, the charge distribution is expected to be approximately uniform 
along the plates, with deviations arising only near the edges due to fringing 
fields. The numerical results confirm this behaviour: as shown in 
Figure~\ref{fig:charge_density}, the charge density is nearly constant over 
the interior of each plate, while a slight variation is observed at the 
boundary nodes, consistent with the well-known edge effect.

It should be noted that the values computed at the outermost nodes may be 
less accurate, as the corresponding elements are more affected by the boundary 
conditions and by the coarser approximation of the fringing field region. 
These edge values should therefore be interpreted with care.

\textbf{Capacitance.}
\label{sec:capacitance_remarks}

Both methods are expected to yield a value slightly exceeding the ideal capacitance
\begin{equation}
    C_{\text{ideal}} = \varepsilon_0 \frac{L_x}{L_y},
\end{equation}
due to the contribution of fringing fields, with a discrepancy of about 5--10\% for 
the chosen aspect ratio $L_x/L_y = 10$. Since the energy method integrates over the 
entire domain $\Omega$, it naturally captures this contribution without any additional 
correction. The charge method, on the other hand, recovers the same information 
indirectly through the normal flux on the plate surface, and the two estimates are 
therefore expected to be mutually consistent.
Any discrepancy between the two methods provides a practical indicator of the 
numerical accuracy of the solution: a relative difference below 1\% is considered 
satisfactory for the purposes of this study.

\subsection{Validation Results}

The coupled FEM-BEM solver is tested on the parallel-plate capacitor described above 
($L_x = 0.10$ m, $L_y = 0.01$ m, $L_x/L_y = 10$, $V_1 = 1$ V, $V_2 = 0$ V) using two 
mesh refinement levels: a coarse mesh with $\delta = 0.005$ m and a fine mesh with 
$\delta = 0.002$ m. The reference value for the capacitance per unit length, neglecting 
fringing fields, is $C_{\text{ideal}} = \varepsilon_0\, L_x/L_y = 88.5$ pF/m.

\paragraph{Electric potential.}
Figures~\ref{fig:potential_comparison} show the electric potential distribution for the 
two mesh levels. In both cases the potential is monotone between the plates, ranging from 
$V_2 = 0$ V at the lower plate to $V_1 = 1$ V at the upper plate, and decays smoothly 
toward $\Gamma$ with no visible oscillations. The equipotential lines are horizontal and 
uniformly spaced in the central region between the plates, confirming the expected 
uniform-field behaviour, and curve regularly outward near the plate edges — the 
characteristic signature of fringing fields. The coarse mesh ($\delta = 0.005$ m) 
reproduces the overall potential distribution correctly but with noticeably lower 
resolution: the equipotential lines appear jagged and the field curvature at the edges 
is poorly resolved. The fine mesh ($\delta = 0.002$ m) produces a significantly smoother 
solution, capturing both the uniform region and the fringing zone with greater fidelity. 
This improvement is expected, as the finer discretization provides more degrees of 
freedom in the regions where the potential gradient changes rapidly.


\begin{figure}[!htbp]
    \centering

    % ---- Row 1: potential maps ----
    \begin{subfigure}{0.48\textwidth}
        \centering
        \includegraphics[width=\textwidth]{Documentation/runs/mesh_0020/margin_10/png/potential_map.png}
    \end{subfigure}
    \hfill
    \begin{subfigure}{0.48\textwidth}
        \centering
        \includegraphics[width=\textwidth]{Documentation/runs/mesh_0020/margin_10/png/potential_contour.png}
    \end{subfigure}

    \caption{Electric potential distribution (left) and equipotential lines (right)}
    \label{fig:potential_comparison}
\end{figure}
% ------------------------------------------------------------------

\begin{figure}[!htbp]
    \centering

    % --- Prima riga ---
    \begin{subfigure}{0.48\textwidth}
        \centering
        %\fbox{\rule{0pt}{0.7\textwidth}\rule{0.9\textwidth}{0pt}}
        \includegraphics[width=\textwidth]{Documentation/runs/mesh_0050/margin_10/png/Efield_map.png}
        \caption*{$E_{field} \delta = 0.005$ m}
    \end{subfigure}
    \hfill
    \begin{subfigure}{0.48\textwidth}
        \centering
        %\fbox{\rule{0pt}{0.7\textwidth}\rule{0.9\textwidth}{0pt}}
        \includegraphics[width=\textwidth]{Documentation/runs/mesh_0020/margin_10/png/Efield_map.png}
        \caption*{$E_{field} \delta = 0.002$ m}
    \end{subfigure}

    \medskip

    % --- Seconda riga ---
    \begin{subfigure}{0.48\textwidth}
        \centering
        \includegraphics[width=\textwidth]{Documentation/runs/mesh_0020/margin_10/png/Efield_bad_elements.png}
        \caption*{Element with a bad E}
    \end{subfigure}
    \hfill
    \begin{subfigure}{0.48\textwidth}
        \centering
        \includegraphics[width=\textwidth]{Documentation/runs/mesh_0020/margin_10/png/Efield_histogram.png}
        \caption*{Distribution of E intensity}
    \end{subfigure}

    \caption{%
    Electric field intensity and direction (top row).
    Elements with $|\mathbf{E}| > 105$~V/m ($>5\%$ of $|\mathbf{E}|$)highlighted in red  and distribution of the element(bottom row).
    }
    \label{fig:efield_comparison}
\end{figure}\textbf{}


\paragraph{Electric field.}
Figures~\ref{fig:efield_comparison} show the electric field 
magnitude and direction for the two mesh levels. In both cases the field is directed 
vertically between the plates, as expected, with magnitude close to the theoretical 
value $|\mathbf{E}| = (V_1 - V_2)/L_y = 100$ V/m in the central region.

The coarse mesh ($\delta = 0.005$ m) resolves the field poorly: individual elements 
show significant variability, with values alternating above and below the expected 
uniform level, reflecting the inability of the coarse discretization to capture the 
smooth field gradient accurately. No element in this case exceeds the 
threshold $E_{\text{threshold}} = 1.05 \times E_{\text{theoretical}} = 105$ V/m, meaning that 
the singularity at the plate edges is not resolved at all at this resolution — the 
mesh is too coarse to capture the field concentration, and the anomalous elements 
are simply smeared into their neighbours.

The fine mesh ($\delta = 0.002$ m) produces a markedly better result: the field is 
visibly uniform across the central region between the plates, and the quiver plot 
confirms the vertical orientation of $\mathbf{E}$ away from the edges. A small number 
of elements near the plate corners exceed the threshold, and their spatial distribution 
matches precisely the physically expected pattern — field concentration at the tips of 
infinitely thin conductors is a well-known singularity of this geometry and is not an 
artefact of the numerical method. The fact that these anomalous elements appear only 
with the finer mesh, and are localised at the corners, confirms that the coupled solver 
is behaving correctly: the finer discretization resolves the singularity rather than 
hiding it.



% ------------------------------------------------------------------
\paragraph{Normal derivative on $\Gamma$.}
Figures~\ref{fig:q_gamma} show the normal derivative $\partial u / \partial n$ on 
$\Gamma$ for the two mesh levels. Recalling that $\partial u / \partial n = \mathbf{E} 
\cdot \mathbf{n}$, this quantity represents the normal component of the electric field 
on the artificial boundary, i.e. the flux per unit length leaving $\Omega_{\text{int}}$ 
through $\Gamma$.

From a physical standpoint, the field lines connecting the two plates are contained 
within $\Omega_{\text{int}}$ and do not cross $\Gamma$; only the fringing fields, which 
curl around the plate edges and extend outward, contribute to the flux through $\Gamma$. 
Consequently, $\partial u / \partial n$ is expected to be negligible on the lateral sides 
of $\Gamma$, which are far from the plates, and to take small but non-zero values on the 
top and bottom sides, in correspondence with the fringing field region above and below 
the plates.

A fundamental consistency check is provided by Gauss's law: since the two plates carry 
equal and opposite charges, the total enclosed charge is zero and the net flux through 
$\Gamma$ must vanish:
\begin{equation}
    \oint_{\Gamma} \frac{\partial u}{\partial n} \, ds = \sum_i q_i \approx 0.
\end{equation}
The numerical value obtained for the fine mesh ($\delta = 0.002$ m) is $\sum_i q_i = -5.47 \times 10^{-13}$. The value is negligibly small in absolute terms, confirming 
that the coupling condition is satisfied and the net flux is numerically consistent with zero.
\begin{figure}[h]
    \centering
    \includegraphics[width=\textwidth]{Documentation/runs/mesh_0020/margin_10/png/q_gamma.png}
    \caption{Normal derivative $\partial u / \partial n$ on $\Gamma$}
    \label{fig:q_gamma}
\end{figure}

% ------------------------------------------------------------------
\paragraph{Surface charge density and Capacitance.}

Figures~\ref{fig:charge_density} show the projected surface charge density 
on the two plates. The distribution is approximately uniform in the central 
region of the plates and increases toward the edges, consistent with the 
well-known behaviour of a finite parallel-plate capacitor where field 
concentration at the tips produces a higher local charge density. The two 
plates carry distributions of opposite sign, as required by the problem setup.

The total charge per unit length on each plate is computed via 
\eqref{eq:total_charge}, and the results for both mesh resolutions are 
reported in Table~\ref{tab:charge}.

\begin{table}[h!]
    \centering
    \begin{tabular}{lccc}
        \hline
        Mesh & $Q_{\text{top}}$ (C/m) & $Q_{\text{bot}}$ (C/m) & 
        $Q_{\text{top}} + Q_{\text{bot}}$ (C/m) \\
        \hline
        Coarse ($\delta = 0.005$ m) & $8.8500 \times 10^{-11}$ & 
        $-8.8500 \times 10^{-11}$ & $1.29\times10^{-26}$ \\
        Fine ($\delta = 0.002$ m)   & $8.9061 \times 10^{-11}$ & 
        $-8.9061 \times 10^{-11}$ & $5.17\times10^{-26}$ \\
        \hline
    \end{tabular}
    \caption{Total charge per unit length on each plate for two mesh 
    resolutions.}
    \label{tab:charge}
\end{table}

In both cases the charges on the two plates are equal and opposite to within 
numerical precision, confirming global charge neutrality. The residual sum 
$Q_{\text{top}} + Q_{\text{bot}}$ is of the order of $10^{-26}$ C/m, which 
is negligible with respect to the individual plate charges and consistent 
with round-off accumulation.

The capacitance estimates obtained from the two methods are summarised in 
Table~\ref{tab:capacitance}. The charge method yields $C_Q = 8.85 \times 
10^{-11}$ F/m, in exact agreement with the ideal value $C_{\text{ideal}} = 
\varepsilon_0 L_x / L_y$ for the coarse mesh. The energy method, when 
restricted to the inter-plate region, recovers the same value. When integrated 
over the full computational domain, however, it yields $C_W = 1.0668 \times 
10^{-10}$ F/m, exceeding $C_{\text{ideal}}$ by approximately 20.5\%. This 
excess is not a numerical artifact: it reflects the physical contribution of 
the fringing fields, which store a non-negligible amount of electrostatic 
energy in the region outside the gap. This is consistent with the estimate 
given in Section~\ref{sec:capacitance_remarks} for the chosen aspect ratio 
$L_x/L_y = 10$. The relative discrepancy between the FEM-BEM solution and 
the reference FEM result is $0.69\%$, confirming the accuracy of the coupled 
formulation.

\begin{table}[h!]
    \centering
    \begin{tabular}{lcc}
        \hline
        Method & $C$ (F/m) & Relative error \\
        \hline
        Ideal (analytical)           & $8.85 \times 10^{-11}$ & --- \\
        Charge method                & $8.85 \times 10^{-11}$ & $0.00\%$ \\
        Energy method (gap only)     & $8.85 \times 10^{-11}$ & $0.00\%$ \\
        Energy method (full domain)  & $1.0668 \times 10^{-10}$ & $20.54\%$ \\
        \hline
    \end{tabular}
    \caption{Capacitance per unit length computed with the two numerical 
    methods. The energy method is evaluated both over the inter-plate gap 
    and over the full domain $\Omega$.}
    \label{tab:capacitance}
\end{table}




\begin{figure}[!htbp]
    \centering

    \begin{subfigure}{0.48\textwidth}
        \centering
        \includegraphics[width=\textwidth]{Documentation/runs/mesh_0020/margin_10/png/charge_density.png}
        \caption*{$\epsilon_{0}\cdot E_{n}$ computation}
    \end{subfigure}
    \hfill
    \begin{subfigure}{0.48\textwidth}
        \centering
        \includegraphics[width=\textwidth]{Documentation/runs/mesh_0020/margin_10/png/charge_density.png}
        \caption*{Energy Method}
    \end{subfigure}

    \caption{Projected surface charge density on the plates}
    \label{fig:charge_density}
\end{figure}

\paragraph{Different $L_y$}

To further characterise the behaviour of the capacitor, an additional study is carried out by fixing $L_x$ and varying the plate separation $L_y$. As $L_y$ increases, the ideal capacitance $C_{\text{ideal}} = \varepsilon_0 L_x / L_y$ decreases, while the relative contribution of fringing fields is expected to grow. 

From a numerical standpoint, the energy-based method computed only within the capacitor region closely follows the theoretical trend of the capacitance, showing excellent agreement with $C_{\text{ideal}}$. On the other hand, the charge-based method slightly overestimates the theoretical value, although the discrepancy remains limited.

In contrast, the energy method evaluated over the entire computational domain progressively deviates from the theoretical capacitance as the domain size increases. However, this behaviour is consistent with the underlying physics: when the plates are not infinite, part of the electric energy is stored outside the capacitor due to fringing fields. As a result, including the external region in the energy calculation naturally leads to higher estimated capacitance values.

\begin{figure}[h]
    \centering
    \includegraphics[width=\textwidth]{Documentation/runs/mesh_0050/margin_10/comparison_capacitance.png}
    \caption{Capacitance vs Ly, fixed Lx}
    \label{fig:q_gamma}
\end{figure}


%% ------------------------------------------------------------
\subsection{Comparison with Pure FEM}

The coupled FEM--BEM formulation is designed to handle the unbounded exterior domain without introducing artificial truncation conditions. To assess the practical advantage of this approach, we compare the coupled solution against a pure FEM solver applied to the same problem on a truncated domain. In the pure FEM setting, the artificial boundary $\Gamma$ is treated as a homogeneous Dirichlet condition $u = 0$, an approximation that becomes more accurate as $\Gamma$ is moved farther from the plates.

\medskip\noindent
As a reference, a pure FEM simulation with a margin of $100\,L_y$ is adopted as the benchmark solution. We report the potential distribution of the pure FEM solution over the full domain (Figure~\ref{fig:fem_pure}), the pointwise difference between the pure FEM and the FEM--BEM solutions over the reduced domain (Figure~\ref{fig:fem_bem_pointwise_comparison}), and the corresponding $L^2$ relative error, computed as described below, amounts to $5.84\%$.
\medskip\noindent
The numerical error is quantified as a percentage $L^2$ relative norm:

\begin{equation}
    \varepsilon_{L^2} =
    \sqrt{
    \frac{\sum_{i \in \mathcal{N}_c}
    \left( u^{\text{FEM}}_i - u^{\text{FEM--BEM}}_i \right)^2}
    {\sum_{i \in \mathcal{N}_c}
    \left( u^{\text{FEM}}_i \right)^2}
    }
    \times 100\%,
\end{equation}

where $\mathcal{N}_c$ denotes the set of nodes in the common region, $u^{\text{FEM}}_i$ is the pure FEM solution, and $u^{\text{FEM--BEM}}_i$ is the coupled solution interpolated at the same locations. Note that the denominator is evaluated using the pure FEM solution, which serves as the reference benchmark. 
\medskip\noindent
The coupled FEM--BEM method required 27.6 s, compared to 3363 s ($\sim$ 56 m) for the pure FEM reference, yielding a 100 times reduction in computational time.

\medskip\noindent
Having established that the coupled FEM--BEM method achieves satisfactory accuracy, we now compare it against pure FEM solutions computed on progressively smaller domains, in order to highlight the efficiency and convenience of the proposed approach.

\begin{figure}[h]
    \centering 
    \includegraphics[width=\textwidth]{Documentation/runs/mesh_0020/margin_10/png/FEM_10_BEM_10.png}
    \caption{Pure FEM $100\,L_y$}
    \label{fig:fem_pure}
\end{figure}


\begin{figure}[h]
    \centering
    \begin{subfigure}{0.48\textwidth}
        \centering
        \includegraphics[width=\textwidth]{Documentation/runs/mesh_0020/margin_10/png/FEM_10_BEM_10.png}
        \caption*{Pure FEM $100\,L_y$ on BEM domain}
    \end{subfigure}
    \hfill
    \begin{subfigure}{0.48\textwidth}
        \centering
        %\fbox{\rule{0pt}{0.7\textwidth}\rule{0.9\textwidth}{0pt}}
        \includegraphics[width=\textwidth]{Documentation/runs/mesh_0020/margin_10/png/u_err_FEM_100_BEM_10.png}
        \caption*{Pointwise comparison}
    \end{subfigure}
    \caption{Comparison between the coupled FEM--BEM method with a $10\,L_{y}$ margin and pure FEM solutions computed on domains with varying margin sizes.}
    \label{fig:fem_bem_pointwise_comparison}
\end{figure}

\medskip\noindent
\textbf{Accuracy as a function of domain size.}
The margin between the capacitor plates and the artificial boundary $\Gamma$ is varied as $L_y$, $10\,L_y$, $20\,L_y$, $50\,L_y$, and $70\,L_y$, where $L_y$ is the vertical spacing between the plates. The mesh step is kept fixed at $\delta = 0.002\,\text{m}$. For each configuration, the $L^2$ relative error on the potential is computed. The FEM--BEM solution, computed on a fixed domain with margin $10\,L_y$ and validated as sufficiently accurate, serves as the reference solution throughout this comparison.

\medskip\noindent
\textbf{Computational time.}
The wall-clock time of each simulation is measured for both the pure FEM and FEM--BEM approaches, allowing us to assess the trade-off between domain size, accuracy, and computational effort. The total time is decomposed into mesh generation time and system solving time; the former is found to be negligible with respect to the latter. Results are reported in Figure~\ref{fig:fem_bem_vs_fem}.

\begin{figure}[h]
    \centering
    \begin{subfigure}{0.48\textwidth}
        \centering
        \includegraphics[width=\textwidth]{Documentation/runs/mesh_0020/margin_10/png/L2_norm.png}
        \caption*{$L^2$ error in percentage}
    \end{subfigure}
    \hfill
    \begin{subfigure}{0.48\textwidth}
        \centering
        \includegraphics[width=\textwidth]{Documentation/runs/mesh_0020/margin_10/png/time.png}
        \caption*{Computational time comparison}
    \end{subfigure}
    \caption{Comparison between the coupled FEM--BEM method with a $10\,L_{y}$ margin and pure FEM solutions computed on domains with varying margin sizes.}
    \label{fig:fem_bem_vs_fem}
\end{figure}


%% ------------------------------------------------------------
\section{Future Implementations}

The present work represents a proof-of-concept implementation of a coupled FEM-BEM approach. Although the error between the coupled method and pure FEM is low and converges as expected, there is room for further improvement. A non-uniform mesh with locally refined elements in the fringing field regions would better capture the rapid variations of the potential near the plate edges. The adoption of higher-order finite elements could additionally increase the accuracy of the solution. From a computational standpoint, reimplementing the solver in C++ would significantly reduce execution times, which is particularly relevant when comparing against large FEM domains in order to more conveniently assess the error and validate the method.

Finally, it would be of considerable interest to extend the analysis to non-symmetric geometries and to fully three-dimensional configurations, which would bring the coupled FEM-BEM framework closer to practical engineering applications.
%% ------------------------------------------------------------
\section{Conclusions}

%% ------------------------------------------------------------
\section{Acknowledgements}

This report is based on Chapter 6 of the book \textit{Topics in Boundary Element Research} by C.A. Brebbia \cite{brebbia1990}. 
The implementation of the code derives, for the FEM part, from the work of Ozlem Ozgun and Mustafa Kuzuoglu \cite{ozgun2018}, distributed under the GNU General Public License, and for the BEM part, from code developed during the course \textit{Computational Electromagnetics} by Prof. Di Rienzo. 
Furthermore, artificial intelligence tools have been used in the preparation of this report. The full project code is available on GitHub: \url{https://github.com/Michelino21/FemBem_Project}.

%% ------------------------------------------------------------

%% ------------------------------------------------------------
\newpage
\bibliographystyle{unsrt}  % numerazione in ordine di citazione
\bibliography{references}  % il tuo file .bib si chiama references.bib
%% ------------------------------------------------------------

\newpage

\begin{appendices}
\section{Introduzione Italiana}

Un classico esempio di accoppiamento FEM/BEM in un problema elettrostatico è il calcolo del campo elettrico generato da un condensatore a facce piane parallele. In una simulazione di questo tipo è spesso utile visualizzare non solo il comportamento del campo all’interno del condensatore, ma anche quello all’esterno. Nei corsi di Fisica 2, il campo esterno viene di norma assunto nullo; in realtà, a causa degli effetti di bordo, esso non è trascurabile e in alcune applicazioni può avere un ruolo rilevante. Questa è la motivazione principale per analizzare in maggiore dettaglio il comportamento del campo elettrico in un problema di questo tipo.

In questo progetto considereremo il caso relativamente semplice di un condensatore a facce piane parallele (si veda la figura di riferimento). Oltre alle cariche presenti sulle armature, non saranno considerate altre sorgenti nello spazio. L’equazione differenziale da risolvere è la prima legge di Maxwell ~\eqref{eq:prima_legge_maxwell_semplificata_A} trattata esprimendo il campo elettrico come gradiente del potenziale, in assenza di densita di carica nello spazio. ~\eqref{eq:prima_legge_maxwell_semplificata_potenziale_A}

\begin{equation}
	\nabla \cdot E = 0
\label{eq:prima_legge_maxwell_semplificata_A}
\end{equation}

\begin{equation}
	\Delta V = 0
\label{eq:prima_legge_maxwell_semplificata_potenziale_A}
\end{equation}

Un primo approccio ben noto consiste nell’utilizzare il metodo degli elementi finiti (FEM). Si definisce un dominio sufficientemente ampio da contenere il condensatore e si impone sul bordo un potenziale nullo, assumendo che a distanza sufficientemente grande il potenziale tenda effettivamente a zero. Questo metodo presenta tuttavia alcuni svantaggi: richiede domini molto estesi e quindi mesh di grandi dimensioni, con conseguente aumento dei tempi di calcolo. Inoltre la finezza della mesh deve essere omogenea dentro e fuori dal condensatore, anche se la qualità richiesta per ottenere un buon risultato non è la stessa nelle due regioni. Calcolare il campo esterno tramite FEM porta quindi a un significativo spreco computazionale e, in ogni caso, il potenziale non può essere rappresentato fino all’infinito, introducendo un’inevitabile approssimazione.

Il metodo BEM, invece, offre una soluzione più efficiente al problema. Riducendo di un ordine la dimensione del dominio, permette di ottenere notevoli vantaggi in termini computazionali. Per sua natura il BEM consente di calcolare con precisione la variabile incognita in qualsiasi punto interno a un dominio, risolvendo il problema solo sul suo bordo. Il BEM risulta applicabile anche a domini molto estesi o non limitati, a patto di conoscere a priori l’andamento asintotico del potenziale e che esso converga almeno come $1/r^2$. In tali casi, il BEM riesce a rappresentare correttamente il comportamento all’infinito senza introdurre approssimazioni artificiali.

Lo scopo di questo lavoro è quindi combinare FEM e BEM per ottenere risultati più accurati a costi computazionali inferiori. L’idea è sfruttare il FEM per determinare con precisione la soluzione all’interno del condensatore e utilizzare il BEM per estenderla verso l’esterno fino all’infinito in modo esatto. Il meccanismo di accoppiamento tra i due metodi sarà descritto in dettaglio nella sezione dedicata.



\end{appendices}

\end{document}